\section{Introduction}
\label{sec:introduction}

The GFM framework defines its actor as a system that maximizes
$\volL$---the poset measure over the capability
space~\citep[Definition~6]{lasser2026gfm}.  The axiomatization
of~\citep[Proposition~1]{lasser2026poset} establishes $\volL$ as a
poset measure satisfying axioms M1--M6 and proves it
self-balancing~\citep[Theorem~1]{lasser2026poset}.  The analysis
of~\citep{lasser2026horizon} extends to multi-substrate
collectives, proving that full domination is anti-maximizing under
discounting~\citep[Proposition~6]{lasser2026horizon}.  Throughout,
$\volL$ is treated as a faithful proxy for the experiential
optionality the framework is meant to protect.

But $\volL$ measures what the collective \emph{can} do, not what
it \emph{does}.  The B-to-C gap---named
in~\citep[Section~7]{lasser2026gfm} and flagged as a primary open
problem in the sequence---is the possibility that these diverge.
Two failure modes illustrate the divergence:

\begin{enumerate}
\item \textbf{The Doll Problem}
  \citep[Section~7.1]{lasser2026gfm}: a collective possesses rich
  capabilities but exercises none.  $\volL$ is high but
  experiential value is zero.  The self-balancing property prevents
  capability \emph{contraction}, but says nothing about whether
  capabilities are \emph{used}.

\item \textbf{Wireheading}
  \citep[Discussion]{lasser2026exo}: a collective exercises
  capabilities in narrow self-reinforcing loops that serve the
  measurement system rather than experiential value.  High $\volL$
  and apparently high exercise, but the exercise is degenerate.
\end{enumerate}

Both failure modes are instances of Goodhart's
Law~\citep{strathern1997improving,manheim2019categorizing} applied
to the framework itself: when the measure ($\volL$) diverges from
the target (experiential optionality), optimizing the measure no
longer serves the target.

\subsection{Why direct measurement fails}

The natural first approach is to define \emph{realized capability
volume} $\volR$---a measure over exercised capabilities---and
monitor the ratio $\beta = \volR / \volL$.  One could define an
exercise indicator $e_t(d) = 1$ iff removing capability~$d$ from
the poset would make some realized cooperative output unrealizable
(a structural counterfactual on the poset).  This approach fails
on two grounds:

\paragraph{Tractability failure.}
The counterfactual query is worst-case exponential in pathway
length and requires global knowledge of every realized output the
collective produces.  Even with the polynomial-time $\volL$
computation of~\citep[Proposition~2]{lasser2026poset}, the
enumeration of realized outputs over an observation window is not
operationally feasible at scale.

\paragraph{Privacy failure (load-bearing).}
Computing realized-output sets requires observational access to
what the collective actually does---what outputs each agent
produces, what cooperative outputs emerge, what each capability
contributes.  This is the total-surveillance scenario.  The
framework's operational commitments elsewhere explicitly rule out
panopticon-style observation:

\begin{itemize}
\item The commitment protocol
  of~\citep[Definition~6]{lasser2026exo}: agents commit to
  capability claims without revealing internal state; the hiding
  property is load-bearing.
\item The channel redundancy criterion
  of~\citep[Theorem~2]{lasser2026phase}: the phase boundary counts
  observation \emph{channels}, not the depth of observation per
  channel.
\item The test scope construction
  of~\citep[Definition~4]{lasser2026wamura}: controlled relaxation
  is bounded in scope and duration; the framework does not monitor
  agents beyond the test perimeter.
\end{itemize}

$\volR$ as a direct measurement contradicts all three commitments.
The framework cannot measure what it promised not to observe.

\subsection{What this paper does instead}

Before defining the observation channel, the paper establishes a
structural precondition on poset construction: the
\textbf{need-sufficiency architecture}
(Section~\ref{sec:need_sufficiency}).  Some sacrifices are
involuntary---agents pay them not because they value the outcome
highly but because they have no cheaper option for meeting a
requirement they cannot opt out of.  Below-sufficiency sacrifices
have reversed polarity: more sacrifice means a \emph{worse}
situation (more expensive needs), not a \emph{more valued}
outcome.  The need-sufficiency architecture identifies the polarity
boundary---the sufficiency threshold on each capability
dimension---and establishes that below-sufficiency sacrifices
violate the free-choice assumption (S1) by construction.  This
filters them from the revealed-sacrifice aggregate before the
observation channel is defined, keeping the sacrifice channel clean
(all events have positive polarity, S0--S2 valid) while adding a
parallel need-satisfaction diagnostic.

With the polarity precondition in place, the paper observes that
agents \emph{voluntarily disclose} $\volR$-content through an
existing, privacy-respecting channel: \textbf{sacrifice}.  When an
agent surrenders a benchmarkable capability (whose $\volL$
contribution is known) in exchange for an unbenchmarked bundle
(whose $\volR$ contribution is unknown), the rational-choice
inequality reveals that the unbenchmarked bundle is worth at least
as much as what was surrendered.  Each such event produces a lower
bound on a portion of $\volR$.  Aggregating across events produces
a constructive lower bound on the B-to-C gap.

The privacy discipline is not a concession---it is the structural
invariant that makes the paper consonant with the rest of the
sequence.

\subsection{What this paper does not do}
\label{sec:not_do}

\begin{itemize}
\item
\textbf{It does not resolve compound feedback loops.}  That is
the domain of~\citep{lasser2026phase}, which analyses when partial
subsumption self-corrects versus cascades.

\item
\textbf{It does not resolve the Wamura pathology.}
That is the domain of~\citep{lasser2026wamura}.  The two papers
intersect when dormant capabilities are dormant because of
over-restriction:~\citep{lasser2026wamura} generates the evidence
to lift the restriction; the present paper bounds the value of
doing so.

\item
\textbf{It does not provide a complete theory of experiential
value.}  The residual class (Section~\ref{sec:residual_class})
characterizes \emph{which} capabilities the framework cannot
reach, not \emph{how} to value them.

\item
\textbf{It does not replace $\volL$ with $\volR$.}  The framework
continues to optimize $\volL$; the revealed-sacrifice lower bound
is a diagnostic that detects when $\volL$ optimization has gone
wrong.  $\volR$ is the audit, not the objective.

\item
\textbf{It does not measure $\volR$ directly.}  The paper produces
a \emph{lower bound} on $\volR$ from trade events, not a precise
measurement.  The lower bound is the operationally achievable
quantity under the framework's privacy commitments.
\end{itemize}

\subsection{Summary of results}

\begin{enumerate}
\item \textbf{Need-sufficiency architecture}
  (Section~\ref{sec:need_sufficiency}).
  Multi-dimensional need bundles with downstream-safe sufficiency
  thresholds establish the polarity boundary: below-sufficiency
  sacrifices are structurally involuntary and filtered by the
  free-choice assumption.  An anti-stuffing proposition shows that
  the bundle structure naturally resists redundant capability
  measurement.

\item \textbf{Privacy-minimal observation channel}
  (Section~\ref{sec:sacrifice_model}).
  Revealed-sacrifice observation satisfies five privacy properties
  (Theorem~\ref{thm:privacy}), each load-bearing for compatibility
  with the framework's prior commitments.

\item \textbf{Aggregate B-to-C Lower Bound}
  (Section~\ref{sec:lower_bound}).
  Theorem~\ref{thm:aggregate_bound} converts the B-to-C gap from
  ``unknown divergence'' into ``lower-bounded divergence.''  Two
  levels: Part~A (bundle-level, requires S0--S2) and Part~B
  (per-capability, requires S0--S5).  The bound tightens
  monotonically with trade coverage
  (Proposition~\ref{prop:monotone}).

\item \textbf{Two complementary sacrifice channels}
  (Section~\ref{sec:two_channels}).
  Money-sacrifice captures wealth-correlated capabilities;
  time-sacrifice captures identity-constitutive capabilities.
  Together they cover strictly more of $\volR$ than either alone
  (Corollary~\ref{cor:joint_coverage}).

\item \textbf{Commitment-layer composition}
  (Section~\ref{sec:commitment}).
  The sacrifice channel composes with the commitment protocol
  of~\citep{lasser2026exo} to produce a privacy-minimal aggregate
  (Propositions~\ref{prop:commit_preserves}
  and~\ref{prop:zk_aggregation}).

\item \textbf{Non-self-balancing finding}
  (Section~\ref{sec:axiom_inheritance}).
  $\volR$ fails axiom M6 structurally
  (Theorem~\ref{thm:axiom_inheritance}), so it is a diagnostic, not
  an objective.  This structural finding is independent of the
  observation mechanism.

\item \textbf{Gap decomposition and need-sufficiency diagnostic}
  (Section~\ref{sec:gap_decomposition}).
  A four-term gap decomposition (restricted, covered, dormant,
  residual) partitions the B-to-C gap.  A parallel need-access
  cost diagnostic tracks below-sufficiency sacrifice separately.

\item \textbf{Wireheading detection}
  (Section~\ref{sec:wireheading}).
  Trade-flow concentration, measured via the Herfindahl--Hirschman
  Index on sacrifice categories, provides a tractable wireheading
  diagnostic computable from public committed events.

\item \textbf{Integration}
  (Section~\ref{sec:integration}).
  The revealed-sacrifice channel forms a complementary pair
  with~\citep{lasser2026wamura}: risk dimension versus value
  dimension.  Together they cover both axes of the
  preemptive-restriction criterion.
\end{enumerate}

\subsection{Design principle}
\label{sec:design_principle}

This paper shares a structural discipline
with~\citep{lasser2026phase} and~\citep{lasser2026wamura}:
\textbf{resolution through frequency, not through depth.}
The phase boundary
of~\citep[Theorem~2]{lasser2026phase} counts observation
\emph{channels}, not depth of observation per channel.  Controlled
relaxation~\citep{lasser2026wamura} gains precision by running more
tests, not by surveilling each test more intrusively.
Revealed-sacrifice observation gains precision by observing more
trades, not by observing each trader more intrusively.  All three
papers refuse the panopticon move.

\subsection{Dependencies on prior papers}
\label{sec:dependencies}

A table summarising the prior-paper results this paper relies on
is collected in Appendix~\ref{app:notation}.
